\section{K\"ahler exact sequences and base change}
\label{app:kahler}

For completeness, we record the universal-property maps used in
\cref{thm:relative-differentials}.  Given \(A\)-algebras \(B,C\) and
\(D=B\otimes_A C\), define
\begin{equation}
 \delta:D\to
 (\Om_{B/A}\otimes_BD)\oplus(\Om_{C/A}\otimes_CD)
\end{equation}
on pure tensors by
\begin{equation}
 \delta(b\otimes c)=
 (db\otimes(1\otimes c),\;dc\otimes(b\otimes1)).
\end{equation}
This is an \(A\)-derivation.  Conversely, restricting the universal derivation
of \(D\) along \(B\to D\) and \(C\to D\) gives maps from the two summands to
\(\Om_{D/A}\).  The maps are inverse because they agree with the universal
derivations on pure tensors.

Applied iteratively to \(D=\cR\), this yields
\cref{eq:relative-differential-decomposition}.  The tensor factors on the
right are essential: \(\Om_{\cR^{(i)}/A}\) is naturally a module over the local
ring, while the direct sum must be compared as a module over the global ring.

\subsection{Presentation for a quotient}

If \(P=k[x_1,\ldots,x_n]\), \(I=(f_1,\ldots,f_m)\), and \(B=P/I\), the conormal
sequence is
\begin{equation}
 I/I^2\longrightarrow
 \Om_{P/k}\otimes_PB
 \longrightarrow\Om_{B/k}\longrightarrow0.
\end{equation}
Under the basis \(dx_1,\ldots,dx_n\), the first map sends the class of \(f_i\)
to
\begin{equation}
 df_i=\sum_j\frac{\partial f_i}{\partial x_j}dx_j.
\end{equation}
Choosing the displayed generators of \(I\) therefore gives the Jacobian
presentation \cref{eq:jacobian-presentation}.  The presentation need not be
minimal, and the first arrow need not be injective.

\subsection{Evaluation}

For a real point \(p:B\to\R\), tensoring with \(\R\) through \(p\) gives
\begin{equation}
 \R^m\xrightarrow{J_F(p)^{\mathsf T}}\R^n
 \longrightarrow \Om_{B/k}\otimes_{B,p}\R\longrightarrow0.
\end{equation}
This operation retains the derivatives evaluated at \(p\) and discards
polynomial coefficients that vanish under evaluation.  It cannot serve as a
procedure for deciding whether any real point \(p\) exists.
